\chapter{What you need}
\begin{table}[H]\centering\small
\begin{tabularx}{\linewidth}{@{}lL@{}}
\toprule
\textbf{Item} & \textbf{Requirement}\\\midrule
Computer & MSX2, MSX2+ or MSX turbo R\\
Memory & 64 KB RAM and 128 KB VRAM\\
Video & The MSX2's own video output; SCREEN~5, 50 or 60~Hz. No additional video cartridge is needed\\
Sound & Internal PSG; optional OPL4 (MoonSound or compatible)\\
Game & 2 MiB ASCII16 MegaROM cartridge\\
Connection & One free slot for the game; another when using MoonSound\\
Display & Monitor or television connected to the MSX video output\\
Controls & Keyboard; optional MSX joystick in port 1 or MSX mouse in port 1 or 2\\\bottomrule
\end{tabularx}\caption{Required equipment.}\end{table}
No tape or disk is needed. To keep your heroes between sessions, write down their save codes (section~\ref{sec:codigos}).

\chapter{Setting up}
\begin{enumerate}
\item Make sure the computer is \textbf{switched off}.
\item Insert the game cartridge firmly into a free slot. If using MoonSound, insert it into another free slot, also with the computer switched off.
\item Connect your mouse or joystick now if you intend to use one (section~\ref{cap:controles}).
\item Connect the display to the MSX video output.
\item Switch the computer on. The game starts automatically.
\end{enumerate}
\aviso{Take care with the cartridge}{Always insert and remove cartridges with the computer \textbf{switched off}. Never insert or remove one while it is on: you could damage both the cartridge and the computer, and lose the current game.}

\chapter{From power-on to the game}
The following screens appear in this order:
\begin{enumerate}
\item \textbf{The DiHalt logo and its musical fanfare}. This sequence lasts about 9.4 seconds with the internal PSG, including its closing pause, or 10 seconds with MoonSound. It advances automatically and cannot be skipped.
\item \textbf{The opening illustration}. Press \key{SPACE}, the joystick fire button or the left mouse button to continue.
\item \textbf{Language selection}. Choose ENGLISH or ESPAÑOL with up/down and confirm with \key{SPACE}, \key{RETURN} or joystick fire. Spanish is selected by default. The intro and game use this language until reset. Hero codes work in either language.
\item \textbf{The animated introduction}, in which Mentor tells Morcar's story. Skip it with \key{SPACE}, \key{ESC} or the right mouse button.
\item \textbf{The main menu}.
\end{enumerate}
\begin{figure}[H]\centering
\includegraphics[width=0.8\linewidth,height=0.6\linewidth]{opening16}
\caption{The opening illustration as displayed by the MSX2, in sixteen colours.}
\end{figure}
\pantalla[0.62]{language_selector_en}{Choose the language before the introduction.}
\subsection{Language}
The cartridge contains both languages. Your choice applies to the entire introduction, menus, messages and all twenty-four adventures. To choose again, restart the game; the language selection is not saved. Write down your hero codes before restarting if you want to keep those characters. The choice of language does not change the rules or the hero-code format.

The screenshots in this manual show the game in English.

\subsection{What is lost when you switch off}
The cartridge does not save data: switching off loses the current game. To preserve your heroes, the game gives each one a \textbf{16-letter code} containing their name, gold, body points, weapons, armour and quest treasures. Write it down before switching off and enter it another day to recover that hero as they were (page~\pageref{sec:codigos}).

\chapter{Controls}\label{cap:controles}
The game is operated with an \textbf{arrow pointer} and a single \textbf{action button}. You can move it with any of these devices, without configuring anything:
\begin{table}[H]\centering\small
\begin{tabularx}{\linewidth}{@{}lLL@{}}
\toprule\textbf{Device} & \textbf{Move pointer} & \textbf{Action}\\\midrule
MSX mouse (port 1 or 2) & Move the mouse & Left button\\
Keyboard & Cursor keys & SPACE\\
Joystick (port 1) & Stick & Fire button\\\bottomrule
\end{tabularx}\caption{Control devices.}\end{table}
\subsection{Mouse}
The game automatically detects an MSX mouse in port 1 or 2 and uses it alongside the keyboard. The left button performs the action; the right button only skips the introduction. A mouse makes the game easier to operate: point at an icon, square or option and click.
\subsection{Keyboard and joystick}
Without a mouse, play works in the same way: the cursor keys, or the joystick in port 1, move the pointer around the screen, and SPACE or fire clicks whatever it points at.
\subsection{Menus}
A line is always highlighted in menus. With keyboard or joystick, \textbf{up} and \textbf{down} move the highlight and SPACE selects it. Without a mouse the arrow is hidden in menus; it returns during play to operate the icons. With a mouse, move the pointer over the option you want and click the left button; the highlight follows it. Clicking outside the options does nothing.
\begin{table}[H]\centering\small
\begin{tabularx}{\linewidth}{@{}lL@{}}
\toprule\textbf{Key} & \textbf{Function}\\\midrule
Cursor keys & Move pointer or menu highlight\\
SPACE & Action: select, click an icon, stop the movement number, or skip the introduction\\
ESC & Skip the introduction\\\bottomrule
\end{tabularx}\caption{Keyboard summary.}\end{table}
